%! Justifying a Claim Based on a Confidence Interval for a Population Proportion — Unit 6, Lesson 6.3 — Warm-Up (Answer Key)
\documentclass[10pt]{article}
\usepackage{apstats-article}
\usepackage{apstats-key}   % pulls in -boxes; do NOT also load -boxes
\newcommand{\TallMath}[1]{$\displaystyle #1\rule[-1.4em]{0pt}{3.2em}$}

\begin{document}

\pageheader{Unit 6, Lesson 6.3}{Warm-Up --- Answer Key}
\namedateperiod

\vspace{0.15in}

A random sample of 360 college students found 198 who said they regularly eat breakfast.

\begin{enumerate}
  \item Verify the three conditions for a one-sample $z$-interval.
        \begin{itemize}
          \item Random: \ans{SRS stated --- met}
          \item Large Counts: \ans{$360(0.55)=198\ge10$; $360(0.45)=162\ge10$ --- met}
          \item 10\% rule: \ans{$360 \le 0.10\times$(all college students) --- met}
        \end{itemize}

  \item Construct a 95\% confidence interval for $p$, the true proportion of all college
        students who regularly eat breakfast. Show all work.

        $\hat p = $ \ans{$198/360 \approx 0.550$}

        \[
          \hat p \pm z^* \sqrt{\frac{\hat p(1-\hat p)}{n}}
          = \ans{0.550} \pm \ans{1.960}\sqrt{\frac{\ans{0.550(0.450)}}{\ans{360}}}
          = (\ans{0.499},\; \ans{0.601})
        \]

  \item Based on your interval, is the value $p = 0.60$ plausible? Is $p = 0.52$ plausible?
        Explain briefly.

        \ansline{$p=0.60$ is barely outside the interval (0.499, 0.601), so it is not plausible at the 95\% confidence level. $p=0.52$ is inside the interval, so it is a plausible value for the true proportion.}
\end{enumerate}

\end{document}
